\documentclass{article}
\usepackage[T1]{fontenc}
\usepackage{lmodern}
\usepackage{graphicx}
\usepackage{amsmath,amssymb}
\usepackage[a4paper,margin=2cm]{geometry}
\usepackage[absolute,overlay]{textpos}
\setlength{\TPHorizModule}{1mm}
\setlength{\TPVertModule}{1mm}
\usepackage[indonesian]{babel}
\usepackage{hyperref}
\setlength{\emergencystretch}{2em}
% unit: Halaman 30

\begin{document}

% === Halaman 30 (python/native) ===

30

• Contoh:  ROT13(ROTATE) = EBGNGR

• Nama “ROT13” berasal dari net.jokes

 (hhtp://groups.google.com/group/net.jokes)  (tahun 1980)

• ROT13 biasanya digunakan di dalam forum online untuk menyandikan 

jawaban teka-teki, kuis,  canda, dsb 

• Enkripsi arsip dua kali dengan ROT13 menghasilkan pesan semula:



P = ROT13(ROT13(P)) 

 sebab        ROT13(ROT13(x)) = ROT26(x) = x

• Jadi dekripsi cukup dilakukan dengan mengenkripsi cipherteks kembali 

dengan ROT13

\end{document}
